\section{Introduction}
\label{sec:intro}

Out-of-focus optics, atmospheric turbulence and motion all blur images, and every
sensor adds noise. Recovering the sharp image is an ill-posed inverse problem.
Inverse filtering divides by the blur's frequency response, which is close to
zero at high frequencies, so the noise explodes (as seen in the Fourier and
denoising parts of the course). Variational methods stabilise the inversion by
adding a regularisation term weighted by a parameter $\lambda$. Two
difficulties remain:
\begin{enumerate}[leftmargin=*,itemsep=1pt]
\item \textbf{Choosing $\lambda$.} Too small amplifies noise, too large
  over-smooths. In practice $\lambda$ is tuned by hand, or by looking at
  the ground truth, which is not available in real use.
\item \textbf{One $\lambda$ for the whole image.} Flat regions would benefit from
  strong smoothing and edges from almost none. A single global weight has to
  compromise between the two. Edge-preserving regularisers such as TV and TGV
  work around this with non-quadratic penalties, but they still have one
  global $\lambda$, and they bring their own artefacts (staircasing for TV).
\end{enumerate}

MULTI~\cite{multi} tackles both at once. It keeps the simplest, quadratic
Laplacian penalty but uses $N$ weights $\lambda_i$, one per pixel. All of them
are set automatically by a closed-form rule, derived from the \emph{uniform
penalty} principle~\cite{upen}, which makes $\lambda_i$ inversely proportional
to the local curvature of the restored image. Only one scalar, $\epsilon$,
remains to be chosen.

\paragraph{Objectives (as approved).}
\begin{itemize}[leftmargin=*,itemsep=1pt]
\item Implement MULTI (Algorithm~1 of the paper) together with its Newton Projection
  inner solver, using a BCCB/FFT Hessian approximation.
\item Test it on Gaussian and out-of-focus (radius~5) blur at noise levels
  $0.5\%$, $1\%$ and $2.5\%$, and compare against Tikhonov, TV and TGV using RE,
  ISNR and SSIM, with a visual study of edges and smooth regions.
\item If time permits: sensitivity to the neighbourhood size and to $\epsilon$;
  colour images; generalisation to other datasets (Kodak).
\end{itemize}
All objectives were completed, and the optional ones are covered in
Secs.~\ref{sec:ablation}--\ref{sec:color}.

\paragraph{Our contributions beyond reproduction.}
Following the instructor's advice to be creative in evaluation and insight:
\begin{itemize}[leftmargin=*,itemsep=1pt]
\item An \emph{edge vs.\ flat} error decomposition, line profiles and $\lambda$-maps,
  which show \emph{where} pixel-wise adaptivity helps and where it does not.
\item A dense $\epsilon$-sweep, which shows that the best $\epsilon$ is set mainly by
  the \emph{image content}, not by the noise level.
\item The discovery and analysis of \emph{$\lambda$-collapse}, a semi-convergence of
  Algorithm~1 that the paper does not discuss, and a one-line fix (a noise floor
  in the $\lambda$ rule, in the spirit of the discrepancy principle~\cite{morozov})
  that removes the catastrophic failures without hurting the good cases.
\item MULTI-auto: a variant with no hand-tuned parameter and no ground truth, and
  MULTI-autoS, which sets $\epsilon$ from the curvature scale of a cheap pilot restoration;
  both are evaluated on held-out Kodak images. Plus a channel-wise colour extension.
\end{itemize}
